\section{How robust is a hidden transition?}\label{sec:robustness}

The examples impose exact equalities on the bank delay curves. This section asks what survives when they are perturbed, and shows that the price of anarchy has a second, optimum-side transition. All statements refer to the family of Theorem~\ref{thm:main}, with route gap $\Delta(t)=a\delta^m+O(\delta^{m+1})$ and $\beta=q_F'(0)$.

\subsection{A small change of free-flow cost restores the kink}

\begin{proposition}[Perturbation law]\label{prop:perturbation}
Change only the connector cost, to $h_\varepsilon(x)=1+\varepsilon+x^p$. For every sufficiently small $\varepsilon\ne0$ the nearby expansion occurs at a unique $t_\varepsilon$ with
\begin{equation}\label{eq:thresholdperturbation}
 t_\varepsilon-1=\sgn(\varepsilon)(|\varepsilon|/a)^{1/m}+O(|\varepsilon|^{2/m}).
\end{equation}
Its slope jump is strictly positive:
\begin{equation}\label{eq:exactperturbation}
 J_E(\varepsilon):=E_\varepsilon'(t_\varepsilon^-)-E_\varepsilon'(t_\varepsilon^+)=\frac{t_\varepsilon[\Delta'(t_\varepsilon)]^2}{F'(Nt_\varepsilon)+G'(Nt_\varepsilon)}\sim\frac\beta2m^2a^{2/m}|\varepsilon|^{2-2/m}.
\end{equation}
The jump of $P$ is $J_E(\varepsilon)/S_\varepsilon(t_\varepsilon)$, and $S_\varepsilon(t_\varepsilon)\to S(1)>0$.
\end{proposition}
\begin{proof}
The inactive route gap becomes $\Delta(t)-\varepsilon$. By \eqref{eq:D} and its differentiated expansion, $\Delta$ is strictly increasing near $1$ with a single critical point at $1$, which gives \eqref{eq:thresholdperturbation}. At the shifted transition the connector flow has right derivative $\Delta'(t_\varepsilon)/(F'(Nt_\varepsilon)+G'(Nt_\varepsilon))$, because $h_\varepsilon'(0)=0$, and left derivative zero. Differentiating \eqref{eq:Ephysical} gives the exact formula. Substituting \eqref{eq:thresholdperturbation}, $\Delta'(t)\sim ma(t-1)^{m-1}$ and $F'(N)+G'(N)=2/\beta$ gives the asymptotic one. All bank loads near $t_\varepsilon$ stay positive. The all-positive marginal-cost equilibrium near $\gamma$ persists for small $\varepsilon$, so $S_\varepsilon$ is smooth near $t_\varepsilon$, and the quotient rule applies.
\end{proof}

In the five-link example the threshold moves on the scale $|\varepsilon|^{1/3}$, but the kink grows only like $|\varepsilon|^{4/3}$. Small perturbations therefore restore a kink that is much smaller than the perturbation (Section~\ref{sec:practice} gives numbers).

\subsection{A small change of the demand ratio}

Keep the network fixed and replace the demand ray by $(N(1+\eta)t,Nt)$. With $b_0=N/\beta$ the inactive gap is $\Delta_\eta(1+\delta)=a\delta^m+b_0\eta+O(\delta^{m+1}+|\eta\delta|+\eta^2)$.

\begin{proposition}[Demand-ratio unfolding]\label{prop:demandunfolding}
For every sufficiently small $\eta\ne0$ there is a unique nearby transition. It is an expansion, with
\[
 t_\eta-1=-\sgn(\eta)\Bigl(\frac{b_0|\eta|}{a}\Bigr)^{1/m}+O(|\eta|^{2/m}),\qquad
 E_\eta'(t_\eta^-)-E_\eta'(t_\eta^+)\sim\frac\beta2m^2a^{2/m}b_0^{2-2/m}|\eta|^{2-2/m}.
\]
The jump of $P$ is the jump of $E$ divided by $S_\eta(t_\eta)\to S(1)>0$.
\end{proposition}
\begin{proof}
Put $\epsilon_0=|\eta|^{1/m}$ and $\delta=\epsilon_0y$. Dividing the gap by $\epsilon_0^m$ gives $ay^m+\sgn(\eta)b_0+\epsilon_0\varrho(y,\epsilon_0)$ with $\varrho$ analytic. The limiting polynomial has one real root, which is simple, and the implicit function theorem continues it. Any nearby zero has $|\delta|$ bounded above and below by multiples of $\epsilon_0$, because the first two terms dominate, so compactness in $y$ gives uniqueness. At the root $\Delta_\eta'(t_\eta)\sim ma^{1/m}b_0^{(m-1)/m}|\eta|^{(m-1)/m}>0$. Proposition~\ref{prop:localcontact}, with contact order one and curvature tending to $2/\beta$, gives the jump. The all-positive equilibrium near the multiplier $\gamma$ persists under small changes of the demand ratio, so the optimum remains analytic near $t_\eta$.
\end{proof}

For $m=3$, $p=4$ we have $b_0=8/5$, the threshold shifts by $(125|\eta|/42)^{1/3}$ to leading order, and the slope jump is $C_\eta|\eta|^{4/3}$ with $C_\eta^3=5267275776/1953125$, so $C_\eta\approx13.919$.

\subsection{Transverse demand directions}

Now vary the two demands $(q_A,q_B)$ independently near $(N,N)$, with the costs fixed. Let $E_0(q_A,q_B)=q_AF(q_A)+q_BG(q_B)$, and put
\[
 f_1=F'(N)=G'(N)=1/\beta,\qquad f_2=F''(N)=G''(N)>0,\qquad \chi=f_1/2+Nf_2/4>0.
\]
These equalities hold in the family because $m\ge3$, and $f_2>0$ because the bank throughput curves are strictly concave.

\begin{proposition}[Anisotropy at the high-contact demand]\label{prop:anisotropy}
For $u_A=q_A-N$, $u_B=q_B-N$ and $\rho=(u_A^2+u_B^2)^{1/2}$,
\begin{equation}\label{eq:planarexpansion}
 E(q_A,q_B)-E_0(q_A,q_B)=-\chi(u_A-u_B)_+^2+o(\rho^2).
\end{equation}
So $E$ is Fr\'echet differentiable at $(N,N)$, with gradient $(3+N/\beta,\,2+N/\beta)$, but has no second-order Taylor expansion there. Along every straight demand path $q_A=N+v_0\delta$, $q_B=N+w_0\delta$ with $v_0\ne w_0$, both $E$ and $P$ are $C^1$ but not $C^2$. Their second-derivative jumps, active minus inactive side, are $-2\chi(v_0-w_0)^2$ and that quantity divided by $S(N,N)$.
\end{proposition}
\begin{proof}
The inactive gap is $f_1(u_A-u_B)+O(\rho^2)$. The connector residual has $x$-derivative $-2f_1$ at the base point, and the cost $x^p$ contributes $O(\rho^p)$. So the implicit function theorem gives the positive root $x=(u_A-u_B)_+/2+O(\rho^2+\rho^p)$. At an active equilibrium
\[
 E-E_0=-\{q_AF'(q_A)-q_BG'(q_B)\}x+\tfrac12\{q_AF''(q_A)+q_BG''(q_B)\}x^2+O(x^3).
\]
The two braces equal $(f_1+Nf_2)(u_A-u_B)+O(\rho^2)$ and $2Nf_2+O(\rho)$. Substitution gives \eqref{eq:planarexpansion}. The positive part of a quadratic is not a quadratic form, which proves the statements about Fr\'echet differentiability. On a straight path put $\sigma=\sgn(v_0-w_0)$ and $\tau=\sigma\delta>0$ on the active side. Then $x=\tau U(\tau,\tau^{p-1})$ and $E-E_0=\tau^2V(\tau,\tau^{p-1})$ with $U,V$ analytic and $V(0,0)=-\chi(v_0-w_0)^2$. Lemma~\ref{lem:remainder} gives the claims for $E$. The common-power scaling maps a neighbourhood of $(N,N)$ to a neighbourhood of $\gamma(N,N)$, where every equilibrium edge is loaded, so the optimum is jointly analytic there, which gives the claims for $P$.
\end{proof}

For the binomial family,
\[
 \chi=\frac{p\{2(m+2)^2+(p-1)(m+1)(m+4)\}}{2N(m+2)^3}.
\]
In the five-link example $\chi=67/125$ and $\nabla E(4,4)=(23/5,18/5)$. The high smoothness therefore belongs to the proportional ray, not to the demand point. Along every other straight path through the same demand vector the second derivative jumps. The expansion \eqref{eq:planarexpansion} is pointwise and does not say that $E$ is $C^1$ on a whole neighbourhood: nearby points of the boundary can have first-derivative jumps.

\subsection{The optimum changes its routes first}\label{sec:optimumside}

By \eqref{eq:optflow} the optimum at $t$ uses the equilibrium loads at $\gamma t$, scaled by $1/\gamma$. Its connector therefore activates at $t=1/\gamma$, before the equilibrium connector does. In the five-link example $1/\gamma=5^{-1/4}\approx0.669$, and at $t=1$ the optimum already sends $0.00643$ units over the connector.

\begin{theorem}[Optimum-side transition]\label{thm:optimumtransition}
In the family of Theorem~\ref{thm:main}, let $V$ be the Beckmann value along the equilibrium ray and $V_0$ its analytic continuation without the connector. Then
\begin{equation}\label{eq:Vexpansion}
 V(T)-V_0(T)=-\frac{k_E}{2m}(T-1)_+^{2m}+\text{(derivative-controlled higher order)},
\end{equation}
and $S(t)=V(\gamma t)/\gamma$. Hence $S\in C^{2m-1}\setminus C^{2m}$ at $t=1/\gamma$, with
\[
 S^{(2m)}(1/\gamma^+)-S^{(2m)}(1/\gamma^-)=-(2m-1)!\,k_E\,\gamma^{2m-1}.
\]
If $E$ is $C^{2m}$ near $1/\gamma$, then $P\in C^{2m-1}\setminus C^{2m}$ at $1/\gamma$, with
\begin{equation}\label{eq:Poptjump}
 P^{(2m)}(1/\gamma^+)-P^{(2m)}(1/\gamma^-)=\frac{(2m-1)!\,k_E\,\gamma^{2m-1}E(1/\gamma)}{S(1/\gamma)^2}>0.
\end{equation}
In the five-link example $E$ is analytic near $1/\gamma$, so $S,P\in C^5\setminus C^6$ there and the jump \eqref{eq:Poptjump} is $\approx+89.68$.
\end{theorem}
\begin{proof}
Proposition~\ref{prop:Sregular} gives $V'(T)=E(T)/T$ along the ray, and the restricted problem without the connector gives $V_0'=E_0/T$ with $V_0$ analytic near $1$. Since $V=V_0$ for $T\le1$ near $1$,
\[
 V(T)-V_0(T)=\int_1^T\frac{E(\tau)-E_0(\tau)}{\tau}\dd\tau .
\]
By the proof of Theorem~\ref{thm:main}, $E-E_0=\delta^{2m-1}V_E(\delta,\delta^{m(p-1)})$ on the right. Dividing by $\tau=1+\delta$ and integrating the absolutely convergent series term by term gives $\delta^{2m}\widetilde V(\delta,\delta^{m(p-1)})$ with $\widetilde V$ analytic and $\widetilde V(0,0)=-k_E/(2m)$. This is \eqref{eq:Vexpansion}, and Lemma~\ref{lem:remainder} applies. The formula $S(t)=V(\gamma t)/\gamma$ is \eqref{eq:optcost}. Substituting $T=\gamma t$ gives $S-S_0=-\frac{k_E\gamma^{2m-1}}{2m}(t-1/\gamma)_+^{2m}+\cdots$, which yields the jump of $S^{(2m)}$. If $E$ is $C^{2m}$, the highest-order jump of $P=E/S$ comes only from $1/S$, and $(1/S)^{(2m)}$ jumps by $-S^{(2m)}\text{-jump}/S^2$. For the five-link example, exact interval arithmetic shows that on $[1/\gamma-1/100,1)$ the equilibrium connector is strictly unused and all four bank edges are loaded, so $E$ is analytic near $1/\gamma$. The numerical value uses $E(1/\gamma)\approx11.920$, $S(1/\gamma)=17/5^{1/4}\approx11.369$ and $P(1/\gamma)\approx1.0485$. \status{proof; exact interval check; numerical constants}
\end{proof}

So a planner who watches the price of anarchy sees two hidden transitions as demand grows: the optimum's at $1/\gamma$, with five continuous derivatives in the quartic example, and the equilibrium's at $1$, with four. In CDS Example~4.3 the optimum also changes regime at a point where $S$ stays smooth. Theorem~\ref{thm:optimumtransition} adds the exact order of that smoothness for the present family.
